%==============================================================================
\section{Conclusions and further work}\label{sec:concl}
%==============================================================================

\subsection{What this report establishes}\label{ssec:concl-summary}

One idea organises everything above: the equivalence between the SSP and a walk
through magazine configurations. It gives a structural theory of the standard
heuristic and a family of exact methods, and the same quantity turns out to govern
both.

\paragraph{On the structural side.}
The two-phase heuristic admits two readings, and separating them is a prerequisite for
saying anything correct about it (Section~\ref{ssec:frames}). Within that setting:

\begin{itemize}[topsep=3pt,itemsep=2pt]
  \item The optimum equals the cheapest walk over \emph{all} feasible groupings
  (Theorem~\ref{thm:grouping}), so the heuristic's loss is exactly the price of using
  the fewest groups (Corollary~\ref{cor:gap}).
  \item Zero-gap classes exist for every capacity: two groups, or a tool set barely
  exceeding the magazine, or a small optimum (Corollaries~\ref{cor:zerogap} and
  \ref{cor:smallZ}).
  \item The ratio is bounded for every instance by
  $\min(b,\Kst-1,q)$ (Theorem~\ref{thm:uncond}), which tightens automatically on
  instances with few groups or a small tool excess.
  \item At three groups the walk cost has a closed form
  (Proposition~\ref{prop:hk3}), and the gap obeys a counting bound
  (Proposition~\ref{prop:k3}) that is attained by explicit witnesses. The natural
  guess $\mathrm{gap}\le\Kst-2$ is false for the walk value
  (Proposition~\ref{prop:refute}).
  \item For the method as actually specified, positive gaps are rare, always equal to
  one in everything tested, and always require $b\ge4$
  (Observations~\ref{obs:rings}--\ref{obs:census}). Unboundedness survives, but with a
  different witness family from the one the walk model suggests
  (Observation~\ref{obs:unbounded-true}).
  \item The circuit clutter of the grouping system computes $\Kst$
  (Proposition~\ref{prop:chromK}) but does not determine the gap
  (Proposition~\ref{prop:noclutter}), while blocking duality gives a first
  non-integral family for the grouping relaxation
  (Proposition~\ref{prop:oddring}).
  \item A per-changeover setup cost interpolates between the SSP and the JGP, with the
  transition confined to an interval whose width is the gap itself
  (Theorem~\ref{thm:collapse} and Lemma~\ref{lem:lowrho}). This answers a question
  posed by \citet{Burger2015}.
\end{itemize}

\paragraph{On the algorithmic side.}
A branch-and-Benders-cut solver was built whose subproblem is exact and polynomial,
together with position-indexed formulations carrying a fixed polynomial row set, among
which the position--transition model is a configuration relaxation that provably
improves on the coverage bound (Proposition~\ref{prop:ptf}).

\paragraph{On the empirical side.}
A campaign over $1{,}301$ instances and eleven configurations, under a single
one-hour limit, establishes that the coverage bound decides the outcome: $94\%$ of
bound-tight instances are closed against $49\%$ of bound-loose ones. Where the solver
fails, its incumbent is already optimal $71\%$ of the time. Fractional Benders cuts,
tested properly for the first time after a defect was corrected, are generated in
enormous numbers and raise the bound on three instances out of $1{,}238$. Of four
strengthenings only the one improving the incumbent helps. The multicommodity model of
\citet{daSilva2024} is reproduced as the strongest single method.

The conclusion those measurements support is that a configuration-based exact method
for the SSP is limited by the strength of its bound and not by its implementation.
This is a negative result about a family of methods, and it is a useful one, because it
says where effort should not go.

\subsection{Further work}\label{ssec:further}

\paragraph{The exact worst case of the heuristic.}
Section~\ref{ssec:law} settles the extremal behaviour of the walk value at three
groups. For the value the method returns, the question is open and is now the more
interesting one, because Section~\ref{ssec:ktnsgap} shows that the two differ
substantially. Two concrete questions: is the heuristic gap ever larger than one on a
connected instance, and does $\mathrm{gap}\le\Kst-2$ hold for it? Directed searches
across $b\in\rset{3,4,5}$ found no counterexample to either.
Problem~\ref{prob:ktns-closes} is the structural question underneath both.

\paragraph{The polyhedral questions.}
Two questions from Section~\ref{ssec:clutter} drive it: whether odd-ring-like
structure is the only obstruction to integrality of the grouping covering relaxation,
and which facets of the position--transition polytope explain where its bound exceeds
$q$. A first scan of the $b=3$ family with two tools per job finds the covering
relaxation non-integral on about a quarter of it, and non-integrality occurring away
from rings, so the obstruction landscape is richer than the odd rings alone.

\paragraph{A bound that reaches into the loose half.}
This is the direction the campaign points to most strongly. The pairwise information of
\citet{Tang1988} is close to exhausted: solving the pairwise bound to optimality, as a
travelling-salesman path over the weights $w_{ij}$ --- an idea already present in their
paper --- lifts the root value on about a quarter of the bound-loose instances but
reaches the optimum on under one per cent of them. The natural next family is the
higher-order form: window inequalities over sets $S$ of consecutive jobs, charging
$\lvert\bigcup_{j\in S}T_j\rvert-b$ when the jobs of $S$ run contiguously, which
generalises the pairwise and triplet rows in one family.

\paragraph{Completing the branch-and-price measurement.}
The runs of PCF$'$ and PTF under the uniform protocol of Section~\ref{ssec:bench} are
in progress. They will answer directly whether PTF's stronger relaxation translates
into solved instances or only into a better root bound.

\paragraph{Non-uniform tool costs.}
Section~\ref{ssec:bbc} notes that the Benders subproblem is the loading linear program
rather than the KTNS rule, and that \citet{Crama1994} show the linear program remains
exact when tools have different setup times while the rule does not. The solver
therefore extends to that regime unchanged. Testing it there would broaden the scope of
the method beyond the uniform problem that the whole SSP literature assumes.

\paragraph{Beyond the SSP.}
Nothing in the position--transition construction is specific to tool switching except
the transition metric. Position-transition columns with per-element consistency rows
apply to any generalised travelling-salesman problem with overlapping clusters
\citep{Pop2024}. Whether the relaxation strength observed here persists in that
generality is open.

\paragraph{Learning-augmented methods.}
Machine learning for the exact side is an active area, covering cut selection
\citep{Tang2020,DezaKhalil2023} and column and branching selection
\citep{Bengio2021,Gasse2019,Morabit2021}. Its established role is to \emph{choose}
among cuts or columns one already knows how to generate, not to make them stronger.
Section~\ref{sec:disc} bounds how much that can help here: what defeats the solver is
a bound gap, and a selector cannot select a strong cut that does not exist. If a
learned component is to help this problem, the more promising target is column
selection for the position-indexed pricers \citep{Morabit2021}, whose coupled-pair
oracle is the measured bottleneck, rather than cut selection for the Benders master.
On the pricing side the corresponding non-learned upgrade is the automatic dual
stabilisation of \citet{Pessoa2018}.
